\section*{\S\,4. Композиция и суперпозиция. Замкнутые множества функций}
\addcontentsline{toc}{section}{\S\,4. Композиция и суперпозиция. Замкнутые множества функций}
\markright{\S\,4. Композиция и суперпозиция. Замкнутые множества функций}
\subsection*{Условия}
1. Пусть $x*y\colon\R^2\to\R$ "--- некоторая ассоциативная операция,
$h(x)\colon\R\to\R$ "--- строго монотонная функция, $h^{-1}(x)$ "---
функция, обратная $h(x)$. Показать, что двухместная операция
$\varphi(x,y)\equiv h^{-1}[h(x)*h(y)]$ ассоциативна.

2. Построить ассоциативную, но не коммутативную двухместную операцию
$x*y\colon\R^2\to\R$.

3. Пусть $h_t(x)=(1-t)e^x+tx$, $h_t^{-1}(x)$ "--- обратная функция
$h_t(x)$. Показать, что функция $\varphi_t(x,y)=h_t^{-1}[h_t(x)h_t(y)]$
комму\-%
тативна и ассоциативна при $0\le t\le1$, а $\varphi_0(x,y)\equiv x+y$,
$\varphi_1(x,y)\equiv xy$.

4. Пусть $A$ "--- функциональное множество, $Q(A)$ "--- наименьшее
замкнутое множество, такое, что $A\subset Q(A)$ ($Q(A)$ "--- замкнутая
оболочка $A$). Показать, что $Q(A)=\M(H)$, где $H$ "--- полная по
отношению к $A$ система функций из $A$. Показать, что при
$H\cap\ov A\ne\emptyset$ $Q(A)\subset\M(H)$. Привести примеры.

5. Пусть $A_i$ ($i=1,\dots,m$) "--- функциональные множества,
$H_i\subset A_i$ "--- соответствующие им полные системы. Показать, что
\[
\bigcup_{i=1}^{i=m}A_i\subset\bigcup_{i=1}^{i=m}\M(H_i)\subset\M\Bigl(\bigcup_{i=1}^{i=m}H_i\Bigr).
\]

6. Функция $\varphi(x_1,x_2)$ определяет ассоциативную двухместную
операцию, если $\varphi(x_1,\varphi(x_2,x_3))\equiv\varphi(\varphi(x_1,x_2),x_3)$.
Аналогично функция $\psi(x_1,x_2,x_3)$ определяет ассоциативную
трехместную операцию, если
\begin{multline*}
\psi(\psi(x_1,x_2,x_3),x_4,x_5)\equiv\psi(x_1,\psi(x_2,x_3,x_4),x_5)\equiv{}\\
\equiv\psi(x_1,x_2,\psi(x_3,x_4,x_5)).
\end{multline*}
Привести пример трехместной ассоциативной операции, не прибегая
к конструкции $\psi(x_1,x_2,x_3)\equiv\varphi(x_1,\varphi(x_2,x_3))$, где
$\varphi(x_1,x_2)$ "--- двухместная ассоциативная операция. Обобщить эти
рассуждения для $n$-местных операций ($n>3$).

7. Построить базисы для множеств целых рациональных функций с целыми,
рациональными и действительными коэффициентами.

\subsection*{Решения}

\begin{solution}[1 (\S\,4)]
\textbf{Проба.} В тексте параграфа среди ассоциативных операций названа
$\varphi_1(x,y)=x+y+\alpha xy$. Посмотрим, почему она ассоциативна. При
$\alpha\ne0$ положим $h(x)=1+\alpha x$; это строго монотонная функция,
отображающая $\R$ на $\R$, и $h^{-1}(u)=(u-1)/\alpha$. Тогда
\[
h(x)\,h(y)=(1+\alpha x)(1+\alpha y)=1+\alpha(x+y+\alpha xy)=h\bigl(\varphi_1(x,y)\bigr),
\]
т.\,е. $\varphi_1(x,y)=h^{-1}[h(x)\,h(y)]$: операция $\varphi_1$ "--- это
обычное умножение, <<пересчитанное>> через $h$. Проверим ассоциативность
прямо: $\varphi_1(\varphi_1(x,y),z)=x+y+\alpha xy+z+\alpha(x+y+\alpha xy)z
=x+y+z+\alpha(xy+xz+yz)+\alpha^2xyz$, и это выражение симметрично
относительно $x$, $y$, $z$, поэтому совпадает с $\varphi_1(x,\varphi_1(y,z))$.
Через $h$ то же видно без выкладки:
$h(\varphi_1(\varphi_1(x,y),z))=h(x)h(y)h(z)=h(\varphi_1(x,\varphi_1(y,z)))$.
Именно это рассуждение и нужно провести в общем случае.

\textbf{Общий случай.} Строго монотонная функция $h$ принимает разные
значения в разных точках, поэтому на множестве ее значений $h(\R)$
определена обратная функция $h^{-1}$, и
\[
h^{-1}(h(u))=u\quad(u\in\R),\qquad h(h^{-1}(v))=v\quad(v\in h(\R)).
\]
Чтобы выражение $h^{-1}[h(x)*h(y)]$ имело смысл при всех $x$, $y$, нужно,
чтобы $h(x)*h(y)$ лежало в $h(\R)$; в условии это подразумевается (так
будет, например, если $h$ отображает $\R$ на $\R$). Вычислим обе части
равенства ассоциативности. Во внутреннем выражении $\varphi(x,y)$ стоит
$h^{-1}[\dots]$, а снаружи к нему применяется $h$; они взаимно уничтожаются:
\begin{align*}
\varphi(\varphi(x,y),z)&=h^{-1}\bigl[h\bigl(h^{-1}[h(x)*h(y)]\bigr)*h(z)\bigr]={}\\
&=h^{-1}\bigl[(h(x)*h(y))*h(z)\bigr],\\
\varphi(x,\varphi(y,z))&=h^{-1}\bigl[h(x)*h\bigl(h^{-1}[h(y)*h(z)]\bigr)\bigr]={}\\
&=h^{-1}\bigl[h(x)*(h(y)*h(z))\bigr].
\end{align*}
Так как операция $*$ ассоциативна, $(h(x)*h(y))*h(z)=h(x)*(h(y)*h(z))$,
и правые части совпадают.

\answer Операция $\varphi(x,y)=h^{-1}[h(x)*h(y)]$ ассоциативна: при
вычислении $\varphi(\varphi(x,y),z)$ и $\varphi(x,\varphi(y,z))$ функции
$h$ и $h^{-1}$ взаимно уничтожаются, и обе части равны
$h^{-1}[h(x)*h(y)*h(z)]$.
\end{solution}

\begin{solution}[2 (\S\,4)]
\textbf{Пробы.} Привычные операции $x+y$, $xy$, $\max(x,y)$, $\min(x,y)$
ассоциативны, но коммутативны. Вычитание не коммутативно, но и не
ассоциативно: $(1-2)-3=-4$, а $1-(2-3)=2$; так же деление и возведение в
степень. Нужна операция, которая по-разному обращается с левым и правым
аргументами, но при этом не <<портит>> группировку.

Самый простой способ "--- вовсе не смотреть на правый аргумент:
$x*y=x$. Тогда $(x*y)*z=x*z=x$ и $x*(y*z)=x$ "--- ассоциативна; но
$1*2=1\ne2=2*1$.

Пример, где участвуют оба аргумента: $x*y=x\,|y|$. Ассоциативность:
\[
(x*y)*z=x\,|y|\,|z|,\qquad x*(y*z)=x\,\bigl|y\,|z|\bigr|=x\,|y|\,|z| .
\]
Некоммутативность: $1*(-1)=1$, а $(-1)*1=-1$. Знак результата берется от
левого аргумента, а модуль "--- от обоих.

\answer Например, $x*y=x\,|y|$ (или совсем простая $x*y=x$): она
ассоциативна, так как $(x*y)*z=x*(y*z)=x\,|y|\,|z|$, и не коммутативна:
$1*(-1)=1\ne-1=(-1)*1$.
\end{solution}

\begin{solution}[3 (\S\,4)]
\textbf{Крайние значения.} При $t=0$: $h_0(x)=e^x$, $h_0^{-1}=\ln$, и
$\varphi_0(x,y)=\ln(e^xe^y)=x+y$. При $t=1$: $h_1(x)=x$, и
$\varphi_1(x,y)=xy$. Так что формула действительно соединяет сложение
и умножение; промежуточные $\varphi_t$ "--- это умножение, <<пересчитанное>>
через $h_t$, как в упражнении~1.

\textbf{План.} По упражнению~1 (с $*$ "--- обычным умножением) операция
$h_t^{-1}[h_t(x)h_t(y)]$ ассоциативна, если $h_t$ строго монотонна и
произведение $h_t(x)h_t(y)$ лежит в множестве значений $h_t$. Напомним
вычисление из упражнения~1:
$\varphi_t(\varphi_t(x,y),z)=h_t^{-1}[h_t(x)h_t(y)h_t(z)]=\varphi_t(x,\varphi_t(y,z))$,
потому что $h_t(h_t^{-1}(v))=v$. Остается проверить два условия.

\textbf{Монотонность.} $h_t'(x)=(1-t)e^x+t$. При $0\le t\le1$ оба слагаемых
неотрицательны и не обращаются в нуль одновременно (при $t<1$ первое
положительно, при $t=1$ второе равно $1$). Значит, $h_t'(x)>0$, и $h_t$
строго возрастает.

\textbf{Множество значений.} При $t=0$ это $(0,\infty)$ "--- произведение
двух положительных чисел положительно. При $0<t\le1$ функция $h_t$
непрерывна, $h_t(x)\to-\infty$ при $x\to-\infty$ (слагаемое $tx$ уходит
в $-\infty$, а $(1-t)e^x\to0$) и $h_t(x)\to+\infty$ при $x\to+\infty$;
поэтому $h_t$ принимает все действительные значения, и $h_t^{-1}$ определена
на всей прямой. В обоих случаях $h_t(x)h_t(y)$ лежит в множестве значений.

\textbf{Коммутативность:} $h_t(x)h_t(y)=h_t(y)h_t(x)$, поэтому
$\varphi_t(x,y)=\varphi_t(y,x)$.

(Явная формула для $h_t^{-1}$ при этом не нужна. Численная проверка
коммутативности и ассоциативности при случайных $t\in(0,1)$ "--- в файле
\texttt{artifacts/checks/associativity.py}.)

\answer При $0\le t\le1$ функция $h_t$ строго возрастает, а произведение
$h_t(x)h_t(y)$ лежит в множестве ее значений, поэтому по упражнению~1
операция $\varphi_t(x,y)=h_t^{-1}[h_t(x)h_t(y)]$ ассоциативна; она
коммутативна, потому что коммутативно умножение. При $t=0$ получаем
$\ln(e^xe^y)=x+y$, при $t=1$ "--- $xy$.
\end{solution}

\begin{solution}[4 (\S\,4)]
\textbf{Напоминание.} $\M(H)$ "--- множество суперпозиций системы $H$:
отображения, которые получаются из отображений системы $H$ и
тождественного отображения конечным числом применений правил в) (составить
пару $(\varphi,\psi)$) и г) (взять композицию $\varphi\circ\psi$). Множество
$\M_0$ \emph{замкнуто}, если $\M(H)\subset\M_0$ для всякой системы
$H\subset\M_0$. Система $H$ \emph{полна} по отношению к $A$, если
$A\subset\M(H)$. Знак $\subset$ в книге означает нестрогое включение
(допускает равенство).

\textbf{Функции и отображения.} Правило в) меняет число компонент: пара
отображений $\R^n\to\R^m$ и $\R^n\to\R^p$ "--- отображение $\R^n\to\R^{m+p}$.
Например, $(x_1+x_2,\ x_1x_2)\colon\R^2\to\R^2$ "--- суперпозиция системы
$H_0$ (см. ниже), но это не многочлен, а пара многочленов. Поэтому множество
числовых функций (непрерывных, многочленов и~т.\,д.) может быть замкнутым,
только если к нему отнесены и векторные отображения, все компоненты которых
ему принадлежат: иначе уже правило в) выводило бы из него. Только при таком
понимании верны примеры 1)--3) замкнутых множеств из текста \S\,4 (в
примерах 4) и 5) множества отображений $\mathfrak X^n\to\mathfrak X^m$ названы
прямо); так будем понимать их и мы. В частности, $C^R$ "--- это \emph{целые
рациональные отображения} $\R^n\to\R^m$ (при любых $n$ и $m$), т.\,е.
отображения, каждая компонента которых "--- многочлен от $x_1,\dots,x_n$.
Слово <<многочлен>> ниже означает многочлен от любого числа переменных
(например, $x_1x_2+x_3$), а при $m=1$ целые рациональные отображения "---
это сами многочлены.

\textbf{Пример для ориентира.} Пусть $A=C^R$ "--- целые рациональные
отображения, $H=H_0=\{x_1+x_2,\ x_1x_2,\ a\ \forall a\in\R\}$. Система $H_0$
состоит из многочленов и полна по отношению к $C^R$: всякий многочлен "---
суперпозиция функций из $H_0$ (пример~1 текста), а отображение, компоненты
которого "--- многочлены, собирается из них правилом в). Множество $C^R$ само
замкнуто (пример~2 замкнутых множеств), так что наименьшее замкнутое
множество, содержащее $C^R$, "--- это $C^R$. Утверждение задачи здесь
говорит: $Q(C^R)=\M(H_0)=C^R$. Так и есть: суперпозиции системы $H_0$ "---
это в точности отображения, все компоненты которых "--- многочлены.

Доказательство разобьем на три шага.

\textbf{1. Замкнутая оболочка существует.} Множество всех правильных
отображений $\mathfrak F(\mathfrak X)$ замкнуто и содержит $A$, так что
замкнутые множества, содержащие $A$, есть. Их пересечение замкнуто
(доказательство теоремы~1 дословно годится для любого числа множеств: если
$H$ лежит в каждом из них, то в каждом лежит и $\M(H)$) и содержится в
любом из них. Это и есть $Q(A)$.

\textbf{2. Множество $\M(H)$ замкнуто.} Пусть $K\subset\M(H)$ "--- любая
система; покажем, что $\M(K)\subset\M(H)$. Для элемента $\M(K)$ назовем
его \emph{рангом} наименьшее число применений правил в), г), с помощью
которого он получается из элементов $K$ и тождественного отображения.
Докажем индукцией по рангу $r$ утверждение
\begin{keystmt}
$P(r)$: каждый элемент $\M(K)$ ранга $r$ лежит в $\M(H)$.
\end{keystmt}
\emph{База}, $r=0$: это элементы $K$ (они лежат в $\M(H)$ по условию) и
тождественное отображение (оно лежит в $\M(H)$ по правилу б)).
\emph{Шаг.} Пусть $P(0),\dots,P(r)$ верны, и элемент имеет ранг $r+1$.
Последнее правило, которым он получен, "--- в) или г), и применено оно
к двум элементам меньшего ранга (их ранги могут быть любыми, не
большими $r$; поэтому нужна полная индукция). По предположению оба лежат
в $\M(H)$, а $\M(H)$ само замкнуто относительно правил в), г). Значит, и
наш элемент лежит в $\M(H)$. Итак, суперпозиции суперпозиций системы $H$
сами являются ее суперпозициями.

\textbf{3. Включения.} (а) Пусть $H$ полна по отношению к $A$. Тогда
$\M(H)$ "--- замкнутое множество (шаг~2), содержащее $A$; так как $Q(A)$
"--- наименьшее такое множество, $Q(A)\subset\M(H)$. Это верно для любой
полной системы, лежит она в $A$ или нет.

(б) Пусть, кроме того, $H$ состоит из функций множества $A$: $H\subset A$.
Тогда $H\subset A\subset Q(A)$, и так как $Q(A)$ замкнуто, $\M(H)\subset Q(A)$.
Вместе с (а) получаем $Q(A)=\M(H)$.

(в) Если же $H\cap\ov A\ne\emptyset$, т.\,е. в $H$ есть функции не из $A$,
то рассуждение (б) не проходит, и остается только включение (а). Оно может
быть как строгим, так и равенством.

\textbf{Примеры.} 1) $H\subset A$, равенство с <<большим>> $Q(A)$: пусть
$A=\{x_1+x_2,\ x_1x_2\}\cup\{a:\ a\in\R\}$ (только эти функции). Система
$H=A$ полна по отношению к $A$, и $Q(A)=\M(A)=C^R$ "--- замкнутая оболочка
может быть гораздо шире самого множества.

2) $H\cap\ov A\ne\emptyset$, строгое включение: $A=C^R$,
$H=H_0\cup\{\sin x\}$. Система полна по отношению к $C^R$, но
$\sin x\in\M(H)$, а $\sin x\notin C^R=Q(A)$ (многочлен, не равный тождественно
нулю, имеет конечное число нулей, а у $\sin x$ их бесконечно много).
Значит, $Q(A)\subset\M(H)$ и $Q(A)\ne\M(H)$.

3) $H\cap\ov A\ne\emptyset$, но равенство: $A$ "--- как в примере 1,
$H=A\cup\{x_1+x_2+x_3\}$. Функция $x_1+x_2+x_3$ не принадлежит $A$, но
является суперпозицией функций из $A$, так что $\M(H)=C^R=Q(A)$.

\answer Если $H\subset A$ и $H$ полна по отношению к $A$, то $Q(A)=\M(H)$:
$\M(H)$ "--- замкнутое множество, содержащее $A$, поэтому $Q(A)\subset\M(H)$,
а $H\subset Q(A)$ и замкнутость $Q(A)$ дают обратное включение. Если же
$H\cap\ov A\ne\emptyset$, остается только $Q(A)\subset\M(H)$; включение может
быть строгим ($A=C^R$, $H=H_0\cup\{\sin x\}$) или равенством
($A=\{x_1+x_2,x_1x_2\}\cup\R$, $H=A\cup\{x_1+x_2+x_3\}$).
\end{solution}

\begin{solution}[5 (\S\,4)]
\textbf{Первое включение.} Система $H_i$ полна по отношению к $A_i$, т.\,е.
$A_i\subset\M(H_i)$; объединяя по $i$, получаем
$\bigcup A_i\subset\bigcup\M(H_i)$.

\textbf{Второе включение} следует из монотонности операции $\M$:
\begin{keystmt}
если $H'\subset H''$, то $\M(H')\subset\M(H'')$.
\end{keystmt}
Действительно, $\M(H'')$ содержит $H'$ и тождественное отображение и
замкнуто относительно правил в), г); тогда каждая суперпозиция системы $H'$
лежит в $\M(H'')$ "--- это та же индукция по числу применений правил,
что в решении упражнения~4 (шаг 2). Применяя монотонность к
$H_i\subset\bigcup_jH_j$, получаем $\M(H_i)\subset\M\bigl(\bigcup_jH_j\bigr)$
для каждого $i$, а значит, и для объединения.

\textbf{Пример: включения могут быть строгими.} Пусть $H_1=A_1=\{x_1+x_2\}$,
$H_2=A_2=\{x_1x_2\}$. Первое включение строгое: например,
$x_1+x_2+x_3\in\M(H_1)$, но не лежит в $A_1\cup A_2$. Для второго заметим, что
компоненты суперпозиций одной функции $x_1+x_2$ "--- это суммы переменных
$x_{i_1}+\dots+x_{i_k}$ (такой вид компонент сохраняется при правилах в), г)),
а компоненты суперпозиций $x_1x_2$ "--- произведения переменных
$x_{i_1}\cdots x_{i_k}$.
Функция $x_1x_2+x_3$ лежит в $\M(H_1\cup H_2)$, но не является ни суммой
переменных (она не линейна), ни произведением: произведение переменных,
равное $1$ и в точке $(1,1,0)$, и в точке $(0,0,1)$, не может содержать ни
$x_3$, ни $x_1$, ни $x_2$.

\answer $\bigcup A_i\subset\bigcup\M(H_i)$, так как $A_i\subset\M(H_i)$ по
определению полноты; $\bigcup\M(H_i)\subset\M(\bigcup H_i)$, так как
$H_i\subset\bigcup H_j$, а $\M$ сохраняет включение. Оба включения могут
быть строгими.
\end{solution}

\begin{solution}[6 (\S\,4)]
\textbf{Пробы.} Нужна функция $\psi(x_1,x_2,x_3)$, у которой все три способа
<<вложить>> $\psi$ в $\psi$ дают одно и то же. Сумма $x_1+x_2+x_3$ и
$\max(x_1,x_2,x_3)$ годятся, но они имеют запрещенный вид
$\varphi(x_1,\varphi(x_2,x_3))$ с $\varphi=+$ или $\varphi=\max$. Медиана
(средняя по величине из трех) не ассоциативна:
$\operatorname{med}(\operatorname{med}(1,2,3),4,5)=\operatorname{med}(2,4,5)=4$,
а $\operatorname{med}(1,\operatorname{med}(2,3,4),5)=\operatorname{med}(1,3,5)=3$.
Попробуем знакопеременную сумму $\psi(x_1,x_2,x_3)=x_1-x_2+x_3$:
\begin{align*}
\psi(\psi(x_1,x_2,x_3),x_4,x_5)&=(x_1-x_2+x_3)-x_4+x_5,\\
\psi(x_1,\psi(x_2,x_3,x_4),x_5)&=x_1-(x_2-x_3+x_4)+x_5,\\
\psi(x_1,x_2,\psi(x_3,x_4,x_5))&=x_1-x_2+(x_3-x_4+x_5).
\end{align*}
Все три равны $x_1-x_2+x_3-x_4+x_5$: операция ассоциативна. Вложенная
<<внутрь>> сумма стоит на месте со знаком $-$ и переворачивает свои знаки,
но и сама начинается с минуса "--- знаки снова чередуются.

\textbf{Почему $x_1-x_2+x_3$ не получается из двухместной ассоциативной
операции.} Допустим, $x_1-x_2+x_3=\varphi(x_1,\varphi(x_2,x_3))$ при всех
$x_1,x_2,x_3$, где $\varphi$ ассоциативна; тогда по ассоциативности также
$x_1-x_2+x_3=\varphi(\varphi(x_1,x_2),x_3)$. Положим $e=\varphi(0,0)$.
При $x_2=x_3=0$ первая запись дает $\varphi(x,e)=x$, а при $x_1=x_2=0$
вторая дает $\varphi(e,z)=z$: число $e$ "--- нейтральный элемент $\varphi$.
Теперь подставим $x_1=e$ в первую запись:
$e-y+z=\varphi(e,\varphi(y,z))=\varphi(y,z)$. Итак, $\varphi(y,z)=z-y+e$
при всех $y$, $z$; но тогда $\varphi(x,e)=e-x+e=2e-x$, а должно быть
$\varphi(x,e)=x$ при всех $x$. Противоречие.

\textbf{Обобщение на $n$-местные операции.} Операция $\psi$ от $n$
аргументов ассоциативна, если при подстановке $\psi$ в любое из $n$ мест
получается одна и та же функция $2n-1$ аргументов:
\begin{multline*}
\psi(\psi(x_1,\dots,x_n),x_{n+1},\dots,x_{2n-1})\equiv{}\\
\equiv\psi(x_1,\psi(x_2,\dots,x_{n+1}),x_{n+2},\dots,x_{2n-1})\equiv\dots\\
\dots\equiv\psi(x_1,\dots,x_{n-1},\psi(x_n,\dots,x_{2n-1})).
\end{multline*}
Запрещенная конструкция "--- $\psi=x_1*x_2*\dots*x_n$ для двухместной
ассоциативной операции $*$.

Знакопеременная сумма обобщается так. Пусть $\zeta$ "--- число, для которого
$\zeta^{n-1}=1$, $\zeta\ne1$, и
\[
\psi(x_1,\dots,x_n)=x_1+\zeta x_2+\zeta^2x_3+\dots+\zeta^{n-1}x_n .
\]
Ассоциативность. Подставим $\psi(x_k,\dots,x_{k+n-1})$ на $k$-е место.
Аргументы $x_1,\dots,x_{k-1}$ сохраняют множители $1,\zeta,\dots,\zeta^{k-2}$.
Внутренняя сумма стоит с множителем $\zeta^{k-1}$, поэтому ее $j$-й аргумент
$x_{k+j-1}$ получает множитель $\zeta^{k-1}\zeta^{j-1}=\zeta^{(k+j-1)-1}$.
Аргументы, стоявшие после $k$-го места, сдвигаются на $n-1$ номеров:
на место $p$ (от $k+1$ до $n$) попадает $x_{p+n-1}$ с множителем
$\zeta^{p-1}=\zeta^{(p+n-1)-1}$ (так как $\zeta^{n-1}=1$). Итак, при любом $k$
каждый $x_m$ получает множитель $\zeta^{m-1}$:
результат $x_1+\zeta x_2+\dots+\zeta^{2n-2}x_{2n-1}$ от $k$ не зависит.

Какие $\zeta$ годятся. Среди действительных чисел $\zeta^{n-1}=1$, $\zeta\ne1$
только при нечетном $n$, $\zeta=-1$: получаем действительную знакопеременную
сумму $x_1-x_2+x_3-\dots+x_n$ (при $n=3$ "--- пример выше). Для любого $n\ge3$
годится комплексное $\zeta=\cos\frac{2\pi}{n-1}+i\sin\frac{2\pi}{n-1}$.
(При четном $n$ действительная знакопеременная сумма не ассоциативна:
например, при $n=4$ подстановка на первое место дает
$x_1-x_2+x_3-x_4-x_5+x_6-x_7$, а на последнее "---
$x_1-x_2+x_3-x_4+x_5-x_6+x_7$.)

Запрещенного вида эти операции не имеют; доказательство то же, что при
$n=3$. Нам понадобится, что $1+\zeta+\dots+\zeta^{n-2}=0$: это сумма
геометрической прогрессии, $(\zeta^{n-1}-1)/(\zeta-1)=0$. Отсюда
$\psi(x,y,\dots,y)=x+y(\zeta+\dots+\zeta^{n-1})=x$ (сумма в скобках
отличается от нулевой заменой $1$ на $\zeta^{n-1}=1$), и
$\psi(y,\dots,y,z)=y(1+\dots+\zeta^{n-2})+\zeta^{n-1}z=z$. Пусть
$\psi=x_1*\dots*x_n$ с ассоциативной $*$, и $e=0*\dots*0$ ($n-1$
множитель). Тогда $x*e=\psi(x,0,\dots,0)=x$ и $e*z=\psi(0,\dots,0,z)=z$:
$e$ "--- нейтральный элемент. Поэтому
$y*z=\psi(y,z,e,\dots,e)=y+\zeta z+e(\zeta^2+\dots+\zeta^{n-1})=y+\zeta z-\zeta e$
(так как $\zeta^2+\dots+\zeta^{n-1}=-\zeta$). Но тогда
$e*z=e+\zeta z-\zeta e$, и равенство $e*z=z$ означает $(\zeta-1)(z-e)=0$ при
всех $z$, что неверно, поскольку $\zeta\ne1$. Противоречие.

(Ассоциативность всех этих операций проверена численно в файле
\texttt{artifacts/checks/associativity.py}.)

\answer Пример: $\psi(x_1,x_2,x_3)=x_1-x_2+x_3$; все три вложения дают
$x_1-x_2+x_3-x_4+x_5$, и эта операция не представляется в виде
$\varphi(x_1,\varphi(x_2,x_3))$ с ассоциативной $\varphi$. Для $n$-местных
операций ассоциативность означает, что подстановка $\psi$ на любое из $n$
мест дает одну и ту же функцию $2n-1$ аргументов; пример "---
$x_1+\zeta x_2+\dots+\zeta^{n-1}x_n$ с $\zeta^{n-1}=1$, $\zeta\ne1$ (при нечетном
$n$ "--- действительная знакопеременная сумма $x_1-x_2+\dots+x_n$); он
тоже не сводится к двухместной ассоциативной операции.
\end{solution}

\begin{solution}[7 (\S\,4)]
\textbf{Что требуется.} Базис в $\M_0$ "--- полная система, которая
перестает быть полной после удаления любой непустой подсистемы. Достаточно
проверять удаление одного элемента: если $H\setminus\{h\}$ неполна при каждом
$h$, то при удалении подсистемы $H_1\ni h$ остается
$H\setminus H_1\subset H\setminus\{h\}$, а по монотонности $\M$ (решение
упражнения~5) $\M(H\setminus H_1)\subset\M(H\setminus\{h\})$, так что и эта
система неполна.

\textbf{Проба.} Система $H_0=\{x_1+x_2,\ x_1x_2,\ a\ \forall a\in\R\}$ из
текста полна для многочленов с действительными коэффициентами, но не
является базисом: например, константу $2$ можно удалить, так как
$2=1+1$ "--- суперпозиция $x_1+x_2$ и константы $1$. Посмотрим, что
действительно необходимо.

Все суперпозиции многочленов "--- целые рациональные отображения (см.
решение упражнения~4): их компоненты "--- многочлены, коэффициенты которых
получаются из коэффициентов исходных функций сложением и умножением
(подстановка многочлена в многочлен и раскрытие скобок). Отсюда три способа
доказать, что функцию \emph{нельзя} получить: если все функции системы
обладают свойством, которое правила в), г) сохраняют у каждой компоненты
(и которым обладает тождественное отображение), то им обладают все
компоненты всех суперпозиций (индукция по числу применений правил, как
в решении упражнения~4).
\begin{itemize}[leftmargin=1.4em,itemsep=1pt]
\item[(а)] <<Многочлен степени не выше 1>> сохраняется: подстановка линейных
  функций в линейную дает линейную.
\item[(б)] <<Одночлен с коэффициентом $\pm1$ или константа $\pm1$>>
  сохраняется: подстановка таких одночленов в такой одночлен дает такой же.
\item[(в)] <<Значение в нуле равно нулю>> сохраняется:
  если $g_i(0,\dots,0)=0$, то $f(g_1,\dots,g_k)$ в нуле равна $f(0,\dots,0)=0$.
\end{itemize}

\textbf{Целые коэффициенты:} $H=\{x_1+x_2,\ x_1x_2,\ -1\}$.

Полнота. Константы: $1=(-1)(-1)$ (функция $x_1x_2$, в которую подставлены
две константы $-1$), $0=1+(-1)$, $n=1+\dots+1$, $-n=(-1)\cdot n$. Переменные
"--- тождественное отображение (с переименованием). Одночлены с целым
коэффициентом $c\,x_1^{a_1}\cdots x_k^{a_k}$ "--- произведения, многочлены "---
суммы одночленов.

Неизбыточность. Без $x_1x_2$ остаются $x_1+x_2$ и $-1$ "--- по (а) все
компоненты суперпозиций линейны, $x_1x_2$ не получить. Без $x_1+x_2$ остаются
$x_1x_2$ и $-1$ "--- по (б) все компоненты суперпозиций "--- одночлены
с коэффициентом $\pm1$ или константы $\pm1$, и $x_1+x_2$ не получить. Без $-1$
остаются $x_1+x_2$ и $x_1x_2$ "--- по (в) все компоненты суперпозиций
обращаются в нуль в нуле, и константу $-1$ не получить. Значит, $H$ "--- базис.

\textbf{Рациональные коэффициенты:}
$H=\{x_1+x_2,\ x_1x_2,\ -1\}\cup\{1/p:\ p\ \text{простое}\}$.

Полнота: как выше получаются все целые числа, а $1/q$ для любого
натурального $q$ "--- произведение чисел $1/p$ по разложению $q$ на простые
множители; затем $m/q=m\cdot(1/q)$ "--- все рациональные константы, и все
многочлены.

Неизбыточность. Без $x_1x_2$ или без $x_1+x_2$ "--- как выше, по (а) и по
(б) (во втором случае коэффициенты одночленов "--- произведения чисел $-1$ и
$1/p$). Без $-1$ все коэффициенты всех суперпозиций неотрицательны
(неотрицательные числа сохраняются при сложении и умножении), и $-1$ не
получить. Без $1/p$ для некоторого простого $p$: все остальные коэффициенты
"--- дроби, в знаменателе которых нет множителя $p$; такие дроби сохраняются
при сложении и умножении, а $1/p$ к ним не относится.

Конечного базиса здесь нет: у конечной системы конечное число
коэффициентов, все их знаменатели делят некоторое $N$, и у всех суперпозиций
знаменатели "--- делители степеней $N$; число $1/p$ для простого $p$, не
делящего $N$, не получить. Поэтому базис неизбежно бесконечен.

\textbf{Действительные коэффициенты.} Здесь ответ зависит от того, как
считать элементы системы. В записи книги $H_0=\{x_1+x_2;\ x_1x_2;\ a,\
\forall a\in\R\}$ константы перечислены одной строкой, как одно семейство.
Если считать все константы одним элементом системы, то $H_0$ "--- базис: без
$x_1x_2$ остаются линейные функции (а), без $x_1+x_2$ "--- одночлены с
действительными коэффициентами (они сохраняются при правилах в), г) так же,
как в (б)), без констант "--- функции, равные нулю в нуле (в).

Если же каждая константа "--- отдельный элемент, то, как показывает проба,
$H_0$ избыточна, и задача становится содержательной. Любая полная
система здесь несчетна: у счетной системы счетное множество коэффициентов,
получаемые из них сложением и умножением числа образуют счетное множество,
а действительных чисел несчетно много. Естественно искать базис вида
$\{x_1+x_2,\ x_1x_2\}\cup C$, где $C$ "--- множество констант. Для этого
достаточно найти множество ненулевых чисел $C$, из которого сложением и
умножением получаются все действительные числа, причем ни одно число из $C$
не получается из остальных. (Свободные члены суперпозиций системы без
константы $c$ получаются сложением и умножением из остальных констант или
равны нулю; поэтому константу $c$ без нее не получить, а $x_1+x_2$ и
$x_1x_2$ необходимы по (б) и (а).) Существует ли такое множество $C$ "--- и
вообще базис в этом случае, "--- мы не выяснили и оставляем вопрос
открытым. (В предисловии автор предупреждает, что часть задач носит
проблемный характер.)

\answer Для целых коэффициентов базис "--- $\{x_1+x_2,\ x_1x_2,\ -1\}$;
для рациональных "--- $\{x_1+x_2,\ x_1x_2,\ -1\}\cup\{1/p:\ p\
\text{простое}\}$ (конечного базиса нет). Для действительных коэффициентов
система $H_0=\{x_1+x_2,\ x_1x_2,\ a\ \forall a\in\R\}$ "--- базис, если
семейство всех констант считать одним элементом. Если каждую константу
считать отдельным элементом, любая полная система несчетна, а базис вида
$\{x_1+x_2,\ x_1x_2\}\cup C$ получится, если найти множество ненулевых чисел
$C$, порождающее все действительные числа сложением и умножением без
избыточности; существует ли оно, осталось открытым.
\end{solution}
